% =============================================================
%  Preamble (book-agnostic).
%  Loaded by main.tex after the document class and geometry.
%  Writers reuse what is defined here. A truly new macro goes in
%  the "Book-specific macros" block at the end and is recorded in
%  book-plan/ledger.md.
% =============================================================

% ---------- Encoding, fonts, language ----------
\usepackage[utf8]{inputenc}
\usepackage[T1]{fontenc}
\usepackage{lmodern}
\usepackage{microtype}
% The book language is set by the scaffold unit from
% BOOK_INSTRUCTIONS.md (e.g. \usepackage[brazilian]{babel}).
% Default: English, no babel.

% ---------- Mathematics ----------
\usepackage{amsmath}
\usepackage{amssymb}
\usepackage{mathtools}
\usepackage{amsthm}

% ---------- Graphics, tables, lists ----------
\usepackage{graphicx}
\usepackage{booktabs}
\usepackage{enumitem}
\usepackage{xcolor}
\usepackage{tikz}
\usetikzlibrary{automata,positioning,arrows.meta}

% ---------- Bibliography (BibTeX + natbib; entries in references.bib) ----------
\usepackage[numbers,sort&compress]{natbib}

% ---------- Code listings ----------
% The ONLY code mechanism of the book. Do not load minted or another
% listings setup.
%   \begin{pythoncode} ... \end{pythoncode}          displayed code
%   \begin{pythoncode}[caption={...},label={lst:x}]   numbered listing
%   \py{expr}                                         inline code
\usepackage{listings}
\definecolor{codebg}{gray}{0.96}
\definecolor{codecomment}{gray}{0.40}
\definecolor{codekeyword}{RGB}{0,70,140}
\definecolor{codestring}{RGB}{120,40,0}
\lstdefinestyle{bookcode}{
  basicstyle=\ttfamily\small,
  keywordstyle=\color{codekeyword}\bfseries,
  commentstyle=\color{codecomment}\itshape,
  stringstyle=\color{codestring},
  backgroundcolor=\color{codebg},
  frame=none,
  columns=fullflexible,
  keepspaces=true,
  showstringspaces=false,
  breaklines=true,
  tabsize=4,
  xleftmargin=0.5em,
  xrightmargin=0.5em,
  aboveskip=0.8\baselineskip,
  belowskip=0.8\baselineskip,
  upquote=true,
}
\lstdefinestyle{pythonstyle}{
  style=bookcode,
  language=Python,
  morekeywords={yield,None,True,False,lambda,nonlocal,match,case},
}
\lstnewenvironment{pythoncode}[1][]{\lstset{style=pythonstyle,#1}}{}
\newcommand{\py}[1]{\lstinline[style=pythonstyle,basicstyle=\ttfamily]{#1}}
\renewcommand{\lstlistingname}{Listing}

% ---------- Theorem-like environments ----------
% Numbered per chapter, one shared counter.
\theoremstyle{plain}
\newtheorem{theorem}{Theorem}[chapter]
\newtheorem{lemma}[theorem]{Lemma}
\newtheorem{proposition}[theorem]{Proposition}
\newtheorem{corollary}[theorem]{Corollary}
\theoremstyle{definition}
\newtheorem{definition}[theorem]{Definition}
\newtheorem{example}[theorem]{Example}
\theoremstyle{remark}
\newtheorem{remark}[theorem]{Remark}
% proof: amsthm's \begin{proof} ... \end{proof}

% ---------- Callout ----------
% \begin{keyidea} ... \end{keyidea}  one short boxed statement of a
% central idea. Use sparingly.
\usepackage[most]{tcolorbox}
\newtcolorbox{keyidea}{
  colback=codebg, colframe=codekeyword, boxrule=0.4pt, arc=1pt,
  left=6pt, right=6pt, top=4pt, bottom=4pt,
  fonttitle=\bfseries, title={Key idea},
}

% ---------- Exercises ----------
% Every substantive chapter ends with:
%   \section{Exercises}\label{sec:<chapter-key>-exercises}
%   \exerciselevel{1}{Foundations}
%   \begin{exercise} ... \end{exercise}
%   \exerciselevel{2}{Application and Reasoning}
%   \exerciselevel{3}{Challenge}
\newcounter{exercise}[chapter]
\renewcommand{\theexercise}{\thechapter.\arabic{exercise}}
\newenvironment{exercise}{%
  \refstepcounter{exercise}%
  \par\medskip\noindent\textbf{Exercise~\theexercise.}\ }{\par}
\newcommand{\exerciselevel}[2]{\subsection*{Level #1 -- #2}}

% ---------- Cross-references ----------
% Label prefixes: part:, chap:, sec:, def:, thm:, lem:, prop:, cor:,
% ex:, rem:, eq:, fig:, tab:, lst:, exr:
\usepackage[hidelinks]{hyperref}
\usepackage[capitalise,noabbrev]{cleveref}
\crefname{exercise}{Exercise}{Exercises}
\crefname{lstlisting}{Listing}{Listings}

% =============================================================
%  Book-specific macros
%  Added by the scaffold unit and by writers. Keep one macro per
%  line with a short comment, and record each one in
%  book-plan/ledger.md.
% =============================================================
